
E. Noether weist dabei noch auf folgendes hin:
Die Richtigkeit von III -- allgemeiner für zyklische Zerfällungskörper beliebigen Grades -- läuft auf den Hasseschen Normensatz\footnote{\label{fn:p38-hasse-normensatz}Angeführt in H, 3.11; bewiesen in: H. Hasse, \emph{Beweis eines Satzes und Widerlegung einer Vermutung über das allgemeine Normenrestsymbol}, Gött. Nachr. 1931.} zurück. Für die Reduktion 3 wird aber nur der Spezialfall des Primzahlgrades, also der ursprüngliche Hilbert-Furtwänglersche Normensatz\footnote{Siehe H. Hasse, Bericht über neuere Untersuchungen und Probleme in der Theorie der algebraischen Zahlkörper II, Jahresber. der D. M.-V., Erg.-Bd. 6 (1930), § 8, sowie die dort angeführte Literatur.} gebraucht. Somit liefert die Reduktion III einen neuen einfachen Beweis des Hasseschen Normensatzes. Dieser Beweis würde so laufen:

Sei \(Z\) ein zyklischer Körper über \(\Omega\) und \(\alpha\) eine Zahl aus \(\Omega\), für die das Normenrestsymbol
\[
\left(\frac{\alpha,Z}{\mathfrak p}\right)=1
\]
ist für jede Primstelle \(\mathfrak p\) von \(\Omega\). Dann zerfällt jede zyklische Algebra
\[
A=(\alpha,Z)
\]
\footnote{\label{fn:p38-hasse-notation}In der Bezeichnung von H, 1. -- Die Angabe des mit \(\alpha\) verknüpften Automorphismus \(S\) von \(Z\) wurde hier als unerheblich unterlassen.} überall [H, 17.7]. Nach der Reduktion 3 folgt also -- unter alleiniger Verwendung des Hilbert-Furtwänglerschen Normensatzes -- \(A\sim\Omega\). Dann ist aber \(\alpha\) Norm eines Elements aus \(Z\) [H, 15.4].

Im Zusammenhang mit dem Normensatz bemerken wir weiter, daß der Satz I als dessen richtige Verallgemeinerung auf höhere (auch nicht-abelsche) Fälle anzusehen ist, während ja die wörtliche Verallgemeinerung, wie H. Hasse zeigte\textsuperscript{\ref{fn:p38-hasse-normensatz}}, nicht allgemein richtig ist.

\begin{center}
\textbf{Folgerungen (H. Hasse).}
\end{center}

\paragraph{5.}
Als nächstliegende Folgerung aus dem Hauptsatz sei angeführt, daß nun die von H. Hasse entwickelte Theorie der zyklischen Algebren über algebraischen Zahlkörpern die Bedeutung einer allgemeinen Struktur- und Invariantentheorie der normalen einfachen Algebren (insbesondere der normalen Divisionsalgebren) über algebraischen Zahlkörpern bekommen hat. Insbesondere ist jetzt allgemein bewiesen:

\paragraph{Satz 1.}
Der Exponent einer normalen einfachen Algebra über einem algebraischen Zahlkörper ist gleich ihrem Index.

Ferner ergibt sich aus Satz I die bemerkenswerte Tatsache:

\paragraph{Satz 2.}
Das Grundideal einer echten normalen Divisionsalgebra über \(\Omega\) ist mindestens durch eine Primstelle von \(\Omega\) teilbar (und sogar mindestens durch zwei).

Dabei ist das Grundideal etwa als die (red.) Norm der (red.) Differente definiert.\footnote{Das läuft im Spezialfall der rationalen Quaternionenalgebren auf den von H. Brandt eingeführten Begriff \glqq{}Grundzahl\grqq{} hinaus (Idealtheorie in Quaternionenalgebren, Math. Ann. 99 (1928)). Da es sich hier nicht um eine Zahl, sondern um ein Ideal handelt, mußte \glqq{}Grundideal\grqq{} gesagt werden; man verwechsle das nicht mit der Dedekindschen Bezeichnung Grundideal und Grundzahl für Differente und Diskriminante, die sich eben aus dem angeführten Grunde für Relativkörper als unbrauchbar erweist. -- Bezüglich der Definition der (red.) Differente siehe die folgende Fußnote. Die (red.) Norm eines Ideals wird von E. Noether ebenfalls \glqq{}primstellenweise\grqq{} definiert, und zwar -- entsprechend der Tatsache, daß für die einzelnen Primstellen jedes Ideal Hauptideal ist -- einfach durch Bildung der (red.) Zahlnormen der Hauptidealbasiszahlen für die einzelnen Primstellen.} Außerdem wird eine unendliche Primstelle dann und nur dann in das Grundideal aufgenommen, wenn sich die Algebra für sie auf die Quaternionenalgebra (und nicht den reellen oder komplexen Zahlkörper) reduziert.

\paragraph{Beweis.}
Nach den Ergebnissen von H. Hasse\footnote{\label{fn:p38-hasse-padisch}H. Hasse, \emph{Über \(p\)-adische Schiefkörper und ihre Bedeutung für die Arithmetik hyperkomplexer Zahlsysteme}, Math. Ann. 104 (1931); siehe dort Satz 42, 59.} geht eine Primstelle \(\mathfrak p\) von \(\Omega\) dann und nur dann im Grundideal auf, wenn der \(\mathfrak p\)-Index von \(1\) verschieden ist. Wären nun alle \(\mathfrak p\)-Indizes gleich \(1\), so zerfiele die Algebra überall, und nach Satz I läge dann keine echte Divisionsalgebra vor. (Daß auch nicht ein \(\mathfrak p\)-Index allein von \(1\) verschieden sein kann, ergibt sich daraus, daß für die Normenrestsymbole, deren Ordnungen die \(\mathfrak p\)-Indizes sind [H, 17.7], das Produkttheorem (Reziprozitätsgesetz) gilt [H, 3.8]).

\paragraph{6.}
Des weiteren ermöglicht Satz I die Bestimmung (arithmetische Kennzeichnung) aller zu einer Algebra \(A\) über \(\Omega\) gehörigen Zerfällungskörper \(K\). In genauer Verallgemeinerung des im Spezialfall zyklischer \(A,K\) von H. Hasse bereits festgestellten Tatbestandes [H, 6, Satz 2] ergibt sich:

\paragraph{Satz 3.}
Für eine Algebra \(A\) über \(\Omega\) ist ein algebraischer Zahlkörper \(K\) über \(\Omega\) dann und nur dann Zerfällungskörper, wenn für die sämtlichen Primteiler \(\mathfrak P_i\) in \(K\) der sämtlichen Primstellen \(\mathfrak p\) von \(\Omega\) jeweils der \(\mathfrak P_i\)-Grad \(n_{\mathfrak P_i}\) von \(K\) ein Multiplum des \(\mathfrak p\)-Index \(m_{\mathfrak p}\) von \(A\) ist.

Dabei ist der \(\mathfrak P_i\)-Grad von \(K\) der Grad des zu \(\mathfrak P_i\) gehörigen \(\mathfrak P_i\)-adischen Körpers \(K_{\mathfrak P_i}\) über dem \(\mathfrak p\)-adischen Körper \(\Omega_{\mathfrak p}\), d. i., wenn
\[
\mathfrak p=\prod_i \mathfrak P_i^{\,e_{\mathfrak P_i}},
        \qquad N_{K\Omega}(\mathfrak P_i)=\mathfrak p^{\,f_{\mathfrak P_i}},
\]
die Zerlegung von \(\mathfrak p\) in \(K\) ist, das Produkt aus dem Grad \(f_{\mathfrak P_i}\) und der Verzweigungsordnung \(e_{\mathfrak P_i}\) von \(\mathfrak P_i\) bzgl. \(\mathfrak p\), \(n_{\mathfrak P_i}=f_{\mathfrak P_i}e_{\mathfrak P_i}\) [H, 4]. Und der \(\mathfrak p\)-Index von \(A\) ist der Index \(m_{\mathfrak p}\) der \(\mathfrak p\)-adischen Erweiterung \(A_{\mathfrak p}\) [H, 6].

\paragraph{Beweis.}
Daß \(K\) Zerfällungskörper für \(A\) ist, d.\,h. \(A_K\sim K\), ist nach Satz I gleichbedeutend damit, daß \(A_{K_{\mathfrak P_i}}\sim K_{\mathfrak P_i}\) für jedes \(\mathfrak P_i\) ist.

Für eine endliche Primstelle \(\mathfrak p\) sei
\[
        A_{\mathfrak p}\sim(\pi,W_{\mathfrak p})
\]
\textsuperscript{\ref{fn:p38-hasse-notation}} die arithmetisch ausgezeichnete zyklische Darstellung von \(A_{\mathfrak p}\) [H, 16; außerdem l. c.\textsuperscript{\ref{fn:p38-hasse-padisch}}, Satz 38], also \(W_{\mathfrak p}\) der unverzweigte Körper vom Grade \(m_{\mathfrak p}\) über \(\Omega_{\mathfrak p}\) und \(\pi\) eine genau durch \(\mathfrak p^1\) teilbare Zahl aus \(\Omega_{\mathfrak p}\). Ferner sei \(K_{\mathfrak P_i}W_{\mathfrak p}\) das Kompositum und
\[
        \Delta_{\mathfrak P_i}=K_{\mathfrak P_i}\cap W_{\mathfrak p}
\]
der Durchschnitt von \(W_{\mathfrak p}\) und \(K_{\mathfrak P_i}\). Da der in \(K_{\mathfrak P_i}\) enthaltene größte unverzweigte Teilkörper vom Grade \(f_{\mathfrak P_i}\) über \(\Omega_{\mathfrak p}\) ist, und da es zu jedem Grad nur einen unverzweigten Körper über \(\Omega_{\mathfrak p}\) gibt, ist \(\Delta_{\mathfrak P_i}\)
\[
        \text{vom Grade }d_{\mathfrak P_i}=(m_{\mathfrak p},f_{\mathfrak P_i})\text{ über }\Omega_{\mathfrak p}.
\]
Daher ist \(K_{\mathfrak P_i}W_{\mathfrak p}\) vom Grade \(m_{\mathfrak p}/d_{\mathfrak P_i}\) über \(K_{\mathfrak P_i}\), und dabei unverzweigt.

Nun ist
\[
        A_{K_{\mathfrak P_i}}=(A_{\mathfrak p})_{K_{\mathfrak P_i}}\sim(\pi,K_{\mathfrak P_i}W_{\mathfrak p})
\]
[H, 15.5]. Demnach ist \(A_{K_{\mathfrak P_i}}\sim K_{\mathfrak P_i}\) gleichbedeutend damit, daß \(\pi\) Norm aus \(K_{\mathfrak P_i}W_{\mathfrak p}\) bzgl. \(K_{\mathfrak P_i}\) ist. Dies ist aber wegen der Unverzweigtheit von \(K_{\mathfrak P_i}W_{\mathfrak p}\) über \(K_{\mathfrak P_i}\) dann und nur dann der Fall, wenn die Ordnungszahl \(e_{\mathfrak P_i}\) von \(\pi\) in \(\mathfrak P_i\) ein Multiplum des Grades \(m_{\mathfrak p}/d_{\mathfrak P_i}\) von \(K_{\mathfrak P_i}W_{\mathfrak p}\) über \(K_{\mathfrak P_i}\) ist, oder also, was wegen
\[
\left(\frac{m_{\mathfrak p}}{d_{\mathfrak P_i}},\frac{f_{\mathfrak P_i}}{d_{\mathfrak P_i}}\right)=1
\]
auf dasselbe hinausläuft, wenn
\[
\frac{f_{\mathfrak P_i}}{d_{\mathfrak P_i}}e_{\mathfrak P_i}=\frac{n_{\mathfrak P_i}}{d_{\mathfrak P_i}}
\]
ein Multiplum von \(m_{\mathfrak p}/d_{\mathfrak P_i}\), d. h. wenn \(n_{\mathfrak P_i}\) ein Multiplum von \(m_{\mathfrak p}\) ist, wie behauptet.

Für den Fall, daß \(\mathfrak p\) eine unendliche Primstelle ist und nicht der triviale Fall \(m_{\mathfrak p}=1\) vorliegt, ist \(m_{\mathfrak p}=2\) und \(A_{\mathfrak p}\) die Quaternionenalgebra über dem reellen Zahlkörper \(\Omega_{\mathfrak p}\). Daß \(A_{K_{\mathfrak P_i}}=(A_{\mathfrak p})_{K_{\mathfrak P_i}}\sim K_{\mathfrak P_i}\) ist, ist dann gleichbedeutend damit, daß \(K_{\mathfrak P_i}\) der komplexe Zahlkörper ist, d. h. mit \(n_{\mathfrak P_i}=f_{\mathfrak P_i}=2=m_{\mathfrak p}\) (\(e_{\mathfrak P_i}\) tritt hier nicht auf), was auch hier wieder die Behauptung ergibt (\(n_{\mathfrak P_i}\) ist wie \(m_{\mathfrak p}\) nur der Werte 1 oder 2 fähig).
